\chapter{Derivable priorities, neutralization, and applicability}
\label{sec:p18-is-dynamic}

This chapter continues the rule-sequence inference system of
Chapter~\ref{sec:p18-is1996}. Its sources are the July 19, 1996 author
version of Prakken and Sartor's dialectical model, Sections 5--8,
and the complete technical companion on defeasible priorities,
especially Sections 5--7 and its proof appendix
\cite{PrakkenSartor1996Dialectical,PrakkenSartor1997Priorities}.
The priority comparison is now extracted from justified arguments.
Its dependence on those arguments must be included in the semantics.
The support-explicit certificates below are an additional construction;
they are not attributed to the historical single-argument game.

\section{Priorities as conclusions}

\begin{hbdefinition}{def:p18-is-dynamic-language}{Dynamic rule theory}
A dynamic theory is \(\Gamma=(S,\mathcal D)\). Arguments, their finite
rule content, \(\operatorname{Con}(A)\), \(\operatorname{Ass}(A)\),
strict extensions, strict closure \(C_S(A)\), and relevant defeasible
frontiers \(R_L(A+T)\) have the meanings in
Chapter~\ref{sec:p18-is1996}. In particular, an argument can have
several conclusions, but has only one producing rule for each.
No global minimality condition is imposed on its content.

Reserve an object-language atom \(x\triangleleft y\), whose arguments
are rule names. It says that \(y\) has priority over \(x\).
The only strict rules containing this predicate are all ground
instances of
\[
\begin{array}{rcl}
x\triangleleft y,\ y\triangleleft z&\to&x\triangleleft z,\\
x\triangleleft y,\ \neg(x\triangleleft z)&\to&
                                      \neg(y\triangleleft z),\\
y\triangleleft z,\ \neg(x\triangleleft z)&\to&
                                      \neg(x\triangleleft y),\\
x\triangleleft y&\to&\neg(y\triangleleft x).
\end{array}
\]
Other strict rules contain no occurrence of this predicate.
Additional positive or negative priority information is supplied by
defeasible rules. A rule name can itself contain names as arguments.
Different ground instances retain different names.
\end{hbdefinition}

\begin{hbdefinition}{def:p18-is-dynamic-comparison}{Comparison relative to support}
Write \(\mathcal A\) for the set of arguments. Attack, coherence and
undercutting \(U\) initially retain their previous definitions,
independent of priorities. A set is conflict-free when there is no
attack between any two of its members, including a member and itself.
For \(X\subseteq\mathcal A\), extract the relation
\[
 E_X=\{(r,s):r\triangleleft s\in\operatorname{Con}(C)
                       \text{ for some }C\in X\}.
\]
This is extraction from conclusions, not an instruction to add the
transitive closure of a relation supplied from elsewhere. Put
\[
 R<_{X}R'\quad\Longleftrightarrow\quad
       \exists r\in R\ \forall s\in R'\ ((r,s)\in E_X).
\]
The existential clause is false for empty \(R\); its universal clause
is true for empty \(R'\).
Define \(\operatorname{Reb}_X(A,B)\) by the previous strict-extension
rebuttal definition, replacing the external frontier comparison by
\(<_X\). Thus some complementary strict-extension conclusions have
frontiers for which the attacking frontier is not \(<_X\) the
attacked frontier. Define
\[
 A D_X B\ \Longleftrightarrow\
 (A=[]\text{ and }B\text{ incoherent})\ \lor\ U(A,B)\
 \lor\bigl(\operatorname{Reb}_X(A,B)\land\neg U(B,A)\bigr).
\]
Finally \(A D_X^sB\) means \(A D_XB\) and not \(B D_XA\).
The subscript is support for comparisons; the arguments being compared
need not belong to that support.
\end{hbdefinition}

\begin{hbtheorem}{thm:p18-is-dynamic-stability}{proved=annex:p18.is-dynamic-stability}{Compatible support and stability of strict defeat}
For conflict-free \(X\), \(E_X\) is asymmetric and irreflexive.
Its set lift \(<_X\) is asymmetric, including the empty-set cases.
The conclusions of this assertion do not require \(E_X\) to be
transitive.
If \(X\subseteq Y\) and both are conflict-free, then
\[
 D_Y\subseteq D_X,\qquad D_X^s\subseteq D_Y^s.
\]
Every attack is oriented by \(D_X\) in at least one direction.
Defeat implies attack except for the designated empty incoherence
case, as in the fixed-priority theory.
\end{hbtheorem}

A conflict-free support set need not itself contain the arguments
combining all its priority conclusions. For example, let its two
members conclude \(r\triangleleft s\) and \(s\triangleleft t\),
respectively, with distinct names and no other defeasible rules.
They are compatible; the extracted relation need not contain
\((r,t)\). The strict transitivity rule licenses a further argument.
It does not insert that further argument into an arbitrary \(X\).
The proofs below use the properties actually established for \(E_X\).

\section{Dynamic justification}

\begin{hbdefinition}{def:p18-is-dynamic-semantics}{The compatible-support operator}
On conflict-free support sets define
\[
 G(X)=\{A:\ \forall B\,[B D_X A\Rightarrow
                       \exists C\in X\ (C D_X^s B)]\}.
\]
Starting with \(J_0=\varnothing\), put
\(J_{\alpha+1}=G(J_\alpha)\) and
\(J_\lambda=\bigcup_{\alpha<\lambda}J_\alpha\) at limit ordinals.
The dynamic justified set \(J\) is the stabilized value.
Justified, overruled and defensible arguments and literals have
the source-specific meanings of Chapter~\ref{sec:p18-is1996};
in particular, overruled arguments are attacked by \(J\).
A theory is finitary here if each argument has finitely many attackers.
This is a restriction on generated arguments, not merely on the
number of displayed rule schemata.
\end{hbdefinition}

\begin{hbtheorem}{thm:p18-is-dynamic-fixed-point}{proved=annex:p18.is-dynamic-fixed-point}{Existence and structural properties}
The operator \(G\) maps conflict-free sets to conflict-free sets and
is monotone on that domain. The displayed iteration is increasing,
stabilizes, and gives the least conflict-free fixed point.
Its value \(J\) is closed under subarguments. If the strict arguments
are inconsistent, \(J=\varnothing\).
If the same displayed formula is evaluated on arbitrary support sets,
every fixed point that fails to contain \(J\) is not conflict-free;
no fixed point is a proper subset of \(J\).
This construction does not treat the conflict-free subsets as a
complete lattice or assert monotonicity on every subset of
\(\mathcal A\).
\end{hbtheorem}

\begin{hbtheorem}{thm:p18-is-dynamic-finitary}{proved=annex:p18.is-dynamic-finitary}{Finite information and finitary approximation}
For fixed finite arguments \(A,B\), each of
\(A D_XB\) and \(A D_X^sB\) depends on only finitely many priority
comparisons between their defeasible rules, even when there are
infinitely many strict extensions.
Consequently, for finitary theories,
\[
 J=\bigcup_{n<\omega}J_n.
\]
Neither statement asserts effective enumeration of all attackers
of an argument in an arbitrary rule language.
\end{hbtheorem}

\section{The source game and an explicit certificate procedure}

\begin{hbdefinition}{def:p18-is-dynamic-game}{Single-argument neutralization game}
The proponent \(P\) opens with the argument \(A\) to be established.
After a proponent argument \(A\), the opponent \(O\) may play any
\(B\) such that \(B D_\varnothing A\).
After this challenge, \(P\) may play an argument \(C\) when either
\[
 C D_{\{C\}}^s B
 \quad\text{or}\quad
 \neg(B D_{\{C\}} A).
\]
The second response is neutralization: \(C\)'s priority conclusions
remove the preceding challenge without necessarily defeating it.
The next opponent move challenges \(C\), again using
\(D_\varnothing\).
The proponent may not repeat argument content on a branch;
the opponent may repeat. A finite proof tree includes every opponent
reply and ends with unchallenged proponent nodes.
This formulation retains the source's dynamic moves, with the
global-minimality restriction removed for the reason proved in
Theorem~\ref{thm:p18-is1996-minimality}.
\end{hbdefinition}

\begin{hbtheorem}{thm:p18-is-dynamic-game-sound}{proved=annex:p18.is-dynamic-game-sound}{Soundness of finite neutralization proofs}
For every dynamic theory, a finite proof tree in this single-argument
game establishes membership in \(J\).
\end{hbtheorem}

\begin{proofboundary}{def:p18-is-dynamic-compression}{Exact scope of the game correspondence}
The 1997 companion's Proposition 6.4 states a completeness result
for finite theories; its proof selects earlier justified support and
combines it into a response argument. The 1996 discussion refers
more broadly to the preservation of finitary results.
Theorem~\ref{thm:p18-is-compression-countermodel} supplies a
finite counterexample even after global minimality is removed.
The precise earlier-rank merger inference fails by
Theorem~\ref{thm:p18-is-compression-rank}: a legal combined
response can be incoherent although its separate supports
are justified. Thus the original finite completeness claim
has no unrestricted finitary extension.

The complete certificate procedure below represents support
separately and proves its own correspondence.
Chapter~\ref{sec:p18-is-dynamic-completeness} then gives a
corrected single-argument-per-move dialogue with sequential
hearings, full-domain soundness and completeness, and finite
winning trees under the finitary hypothesis.
Both additional procedures are explicitly distinguished from
the refuted original single-response game.
\end{proofboundary}

\begin{hbdefinition}{def:p18-is-dynamic-certificate}{Support-explicit certificates}
A certificate for \(A\) is a finite rooted construction. For every
\(B D_\varnothing A\), provide a finite conflict-free set \(X_B\),
and one of the following recorded responses:
\[
 \neg(B D_{X_B} A),\qquad
 \text{or}\qquad C_B\in X_B,\quad C_B D_{X_B}^s B.
\]
For each member of each \(X_B\), provide a separate certificate
below this response. An empty set cannot answer an actual
\(D_\varnothing\) challenge. A leaf claim is permitted exactly
when it has no \(D_\varnothing\)-defeaters. The construction is
well-founded and finite; an argument cannot occur again as a claim
on its own descending support branch.

Equivalently, a game has three position types. At a claim \(A\),
\(O\) chooses a \(D_\varnothing\)-defeater or loses if none exists.
At the resulting challenge \((A,B)\), \(P\) supplies one of the
responses above or loses. At its support position, \(O\) chooses
any member of \(X_B\) and challenges that claim next.
All opponent choices must be covered. An infinite play is not a
winning certificate. This definition specifies its states, moves
and closure conditions without merging several arguments into one.
\end{hbdefinition}

\begin{hbtheorem}{thm:p18-is-dynamic-certificates}{proved=annex:p18.is-dynamic-certificates}{Exact certificate correspondence}
For arbitrary theories every support-explicit certificate is sound.
For finitary theories,
\[
 A\in J\quad\Longleftrightarrow\quad
 A\text{ has a finite support-explicit certificate}.
\]
For a finite set of ground rules these certificates and \(J\) can be
computed by finite search. This is a finite-domain decision result,
not an algorithmic conclusion from finitariness alone.
\end{hbtheorem}

\section{Reversing the fixed-priority roles}

\begin{hbdefinition}{def:p18-is-dynamic-weak-game}{The reversed-role game}
Return in this definition to a fixed defeat relation \(D_\Gamma\)
from Chapter~\ref{sec:p18-is1996}. The proponent opens with \(A\).
The opponent must strictly defeat the preceding proponent argument;
the proponent may use any defeater of the preceding opponent argument.
The opponent cannot repeat argument content on a branch; the
proponent may repeat. A proponent winning tree covers all strict
opponent challenges and ends at proponent nodes with no such reply.
No additional global-minimality restriction is imposed.
This is the source's reversed-role idea at its fixed-comparison
interface. No dynamic reversed-role neutralization rule is
implicitly supplied by it.
\end{hbdefinition}

\begin{hbtheorem}{thm:p18-is-dynamic-weak-game}{proved=annex:p18.is-dynamic-weak-game}{Finite reversed-role characterization}
If the fixed theory has finitely many arguments, the reversed-role
game is determined and
\[
 P\text{ wins from }A
 \quad\Longleftrightarrow\quad
 \text{no member of }J_\Gamma\text{ defeats }A.
\]
A winning argument is therefore not overruled in the source's
attack-based classification. The converse statement about that
classification need not hold because of the empty incoherence
defeat.
\end{hbtheorem}

\begin{hbexample}{def:p18-is-dynamic-residual}{Residual defensibility is not a winning strategy}
Let \(S=\varnothing\), with no priorities, and take just
\(r:\Rightarrow a\) and \(s:\Rightarrow\neg a\).
The only justified content is the empty argument.
The combined argument \([r,s]\) is incoherent. No justified argument
attacks it, so it is in the source's residual defensible class.
But the empty argument strictly defeats it and wins immediately
as the opponent's first move. The annex verifies the entire finite
argument set.
\end{hbexample}

\section{Worked priority dialogues}

\begin{hbexample}{def:p18-is-dynamic-outside-domain}{Why compatibility belongs in the domain}
Take \(r_1:\Rightarrow a\), \(r_2:\Rightarrow b\),
\(r_3:b\Rightarrow\neg a\), \(r_4:\Rightarrow\neg b\),
\(r_5:\Rightarrow r_2\triangleleft r_4\), and
\(r_6:\Rightarrow r_4\triangleleft r_2\).
Let \(X=\{[r_4],[r_5]\}\) and \(Y=X\cup\{[r_6]\}\).
The set \(X\) is conflict-free and \(Y\) is not.
Evaluating the formula for \(G\) even outside its declared domain
gives \([r_1]\in G(X)\) but \([r_1]\notin G(Y)\).
These are membership assertions about the full argument universe,
not equalities omitting strict extensions and combined contents.
\end{hbexample}

\begin{hbexample}{def:p18-is-dynamic-reinstatement}{Derivable priorities restore a weak assumption}
The companion's Example 5.13 starts with
\[
r_0:\Rightarrow a,\quad r_1:a\Rightarrow b,\quad
r_2:\mathord{\sim}b\Rightarrow c,\quad r_3:\Rightarrow\neg a
\]
and adds
\[
r_4:\Rightarrow r_0\triangleleft r_3,\quad
r_5:\Rightarrow r_3\triangleleft r_0,\quad
r_6:\Rightarrow r_5\triangleleft r_4.
\]
The higher priority enters first, the argument using \(r_4\)
is accepted next, and then \([r_3]\) is accepted.
Every argument undercutting \([r_2]\) must use \(r_1\) to conclude
\(b\), hence must also use \(r_0\).
The justified \([r_3]\) strictly defeats each of them.
Consequently \(c\) becomes justified although its weak assumption
has syntactically available counterarguments.
\end{hbexample}

\begin{hbexample}{def:p18-is-dynamic-three-cycle}{Priority axioms expose a conflict between combined supports}
Use the companion's Example 5.14 with just
\[
\begin{array}{lll}
r_4:\Rightarrow r_1\triangleleft r_2,&
r_5:\Rightarrow r_2\triangleleft r_3,&
r_6:\Rightarrow r_3\triangleleft r_1,\\
r_7:\Rightarrow r_6\triangleleft r_4,&
r_8:\Rightarrow r_6\triangleleft r_5.&
\end{array}
\]
The last two defaults and their combination are unattacked.
At the next stage \([r_4,r_5]\) is accepted.
Its strict closure derives \(r_1\triangleleft r_3\), which conflicts
with \(r_6\)'s conclusion. Conversely, strict consequences of
arguments combining \(r_6\) with \(r_4\) or \(r_5\) expose the
other opposed comparisons. Each attacking frontier contains
\(r_6\), whose priority is below both \(r_4,r_5\).
The comparison is between actual argument frontiers; a third
priority completing a cycle is not treated as harmless merely
because its asserted edge differs from the first two.
\end{hbexample}

\begin{hbexample}{def:p18-is-dynamic-neutralization}{A priority response needs its own defense}
Take exactly the defeasible rules
\[
r_1:\Rightarrow a,\quad r_2:\Rightarrow b,\quad
r_3:b\Rightarrow\neg a,\quad r_4:\Rightarrow\neg b,\quad
r_5:\Rightarrow r_2\triangleleft r_4,
\]
and the four strict priority schemata.
A representative branch is
\[
P:[r_1],\quad O:[r_2,r_3],\quad
P:[r_4,r_5],\quad O:[r_2],\quad P:[r_5].
\]
The second proponent move strictly defeats the first challenge
under its own priority. The last move neutralizes the challenge
to \(\neg b\); it does not conclude \(\neg b\) again.
Semantically the priority support is justified at the first stage,
\([r_4]\) at the next, and \([r_1]\) thereafter.
The displayed branch illustrates these moves; the annex covers
all generated challenges for these stage claims.
\end{hbexample}

\begin{hbexample}{def:p18-is-dynamic-circle}{A conclusion cannot supply its own initial priority}
Take
\[
r_1:\Rightarrow a,\quad r_2:\Rightarrow\neg a,\quad
r_3:a\Rightarrow r_2\triangleleft r_1.
\]
If \(r_3\) instead has empty antecedent, its priority argument
is unattacked, and the next stage accepts \([r_1]\).
The displayed conditional antecedent changes this outcome.
The priority argument depends on \(a\), which is precisely the
disputed conclusion it would protect. Apart from the empty content,
no argument is justified. Hence no finite sound neutralization
proof establishes \(a\). With instead just
\(u:\Rightarrow u\triangleleft v\) and
\(v:\Rightarrow v\triangleleft u\), the asymmetry rule makes the
opposed priority arguments conflict, and neither priority is
justified. The source's nested names must not be confused with
permission for circular justification.
\end{hbexample}

\begin{hbexample}{def:p18-is-dynamic-building}{Priorities between two priority-rule instances}
This is a stipulated reconstruction of the source's building example.
Use ordinary rules
\[
R_P(v):\operatorname{Protected}(v)\Rightarrow\neg\operatorname{Modify}(v),
\qquad
R_M(v):\operatorname{Needs}(v)\Rightarrow\operatorname{Modify}(v),
\]
and priority rules
\[
\begin{split}
K(x,y):&\operatorname{Art}(x),\operatorname{Town}(y)
                                  \Rightarrow y\triangleleft x,\\
T(x,y):&\operatorname{Earlier}(x,y)\Rightarrow x\triangleleft y.
\end{split}
\]
Provide unattacked defeasible facts for \(\operatorname{Protected}(v)\),
\(\operatorname{Needs}(v)\), \(\operatorname{Art}(R_P(v))\),
\(\operatorname{Town}(R_M(v))\), and
\(\operatorname{Earlier}(R_P(v),R_M(v))\).
Abbreviate the actual ground rule names by
\[
 K_1=K(R_P(v),R_M(v)),\quad
 T_1=T(R_P(v),R_M(v)),\quad T_2=T(T_1,K_1).
\]
Finally supply the unattacked fact
\(\operatorname{Earlier}(T_1,K_1)\).
There are no other ordinary facts or defeasible rules.

An argument for preservation is challenged by the argument using
\(R_M(v)\). The argument using \(K_1\) supplies
\(R_M(v)\triangleleft R_P(v)\), neutralizing that challenge.
The argument using \(T_1\) challenges this priority with the reverse
comparison. The argument using \(T_2\) supplies
\(T_1\triangleleft K_1\), neutralizing the second challenge.
All the listed factual premises must occur in their corresponding
argument contents. \(T_2\) is a different ground instance from
\(T_1\), not a rule asserting its own superiority.
The unattacked higher priority settles the lower comparison and
then the building conclusion. The annex gives the dependency and
challenge calculation.
\end{hbexample}

\begin{hbexample}{def:p18-is-dynamic-exception}{An exception overrides a ranked rebuttal}
Let \(P_0\) denote a product, and use
\[
\begin{split}
E(x,y,z):\;&\operatorname{Sale}(x,y),\operatorname{EC}(y),
 \operatorname{EC}(z),\mathord{\sim}\operatorname{HealthRisk}(x,z),
 \mathord{\sim}\operatorname{ConsumerHarm}(x,z)
 \Rightarrow\operatorname{Sale}(x,z),\\
L(x):\;&\operatorname{Pasta}(x),\neg\operatorname{HardCorn}(x)
                                      \Rightarrow\neg\operatorname{Sale}(x,I).
\end{split}
\]
For this example take only the displayed ground instances
\(E_0=E(P_0,B,I)\), \(L_0=L(P_0)\), their five required positive
or strongly negative facts, and
\(q:\Rightarrow L_0\triangleleft E_0\).
The five facts are
\(\operatorname{Sale}(P_0,B),\operatorname{EC}(B),\operatorname{EC}(I),
\operatorname{Pasta}(P_0),\neg\operatorname{HardCorn}(P_0)\).
The explicit priority premise realizes the priority assumed in the
source discussion. With no exception argument, \(E_0\)'s conclusion
is justified and \(L_0\)'s opposed conclusion is overruled.
Add \(h:\Rightarrow\operatorname{ConsumerHarm}(P_0,I)\).
It undercuts \(E_0\), independently of \(q\); the local conclusion
is then reinstated. These are computations in a stipulated rule
system, not statements of current legal authority.
\end{hbexample}

\begin{hbtheorem}{thm:p18-is-dynamic-examples}{proved=annex:p18.is-dynamic-examples}{The worked priority outcomes}
The stage and outcome claims in the preceding priority examples hold
for their complete generated argument systems, including strict
extensions and additional compatible argument contents.
\end{hbtheorem}

\section{Two applicability constructions and their limits}

\begin{hbdefinition}{def:p18-is-dynamic-app-encoding}{Applicability in rule premises}
Select a set of original defeasible rules \(r\), and reserve fresh
atoms \(\operatorname{app}(r)\). There are two distinct translations.
The negative-premise translation adds
\(\mathord{\sim}\neg\operatorname{app}(r)\) to each selected rule.
An argument for \(\neg\operatorname{app}(r)\) then undercuts every
use of the translated rule.

The positive-premise translation instead adds
\(\operatorname{app}(r)\) to each selected rule and supplies a fresh
default \(\rho_r:\Rightarrow\operatorname{app}(r)\).
Arguments for inapplicability rebut this default, subject to the
usual priority rules. Priorities comparing its responsible rule
must name \(\rho_r\); an ordering of the original rule \(r\) does
not automatically order \(\rho_r\).
The fresh helper defaults are added after translating the selected
original rules and are not recursively translated.
\end{hbdefinition}

\begin{hbdefinition}{def:p18-is-dynamic-exclusion}{Direct exclusion}
Alternatively extend attack by declaring that \(A\) exclusion-attacks
\(B\) if some defeasible rule \(r\) occurring in \(B\) has
\(\neg\operatorname{app}(r)\in C_S(A)\).
This direction makes the inapplicability argument attack the use
of the rule.
For priority-free exclusion put
\(\operatorname{Ex}_X(A,B)\) equal to this condition.
For priority-sensitive exclusion require, in addition, a witnessing
strict extension \(A+T\) such that
\[
 R_{\neg\operatorname{app}(r)}(A+T)\not<_X\{r\}.
\]
In either proposed variant the corresponding defeat formula is
\[
 A D_X^{\rm ex}B\ \Longleftrightarrow\
 (A=[]\text{ and }B\text{ incoherent})\lor U(A,B)
 \lor\bigl((\operatorname{Reb}_X(A,B)\lor\operatorname{Ex}_X(A,B))
                         \land\neg U(B,A)\bigr).
\]
For this formula coherence and conflict-freedom are measured using
the enlarged attack relation. Exclusion targets defeasible rules
only. Reverse undercutting blocks this exclusion clause.
\end{hbdefinition}

\begin{hbtheorem}{thm:p18-is-dynamic-applicability}{proved=annex:p18.is-dynamic-applicability}{Supported applicability results and distinctions}
Each premise translation is an ordinary instance of the dynamic
core and inherits its proved results.
Priority-free direct exclusion also preserves conflict orientation,
support stability, the compatible-support fixed point, finitary
approximation and the certificate correspondence, using its enlarged
attack relation throughout.
It is not equivalent to the negative-premise translation.
Priority-sensitive direct exclusion, with the displayed enlarged
attack notion and the same operator \(G\), need not map
conflict-free support sets to conflict-free sets. Thus it does not
inherit that fixed-point theorem without another change of semantics.
\end{hbtheorem}

\begin{hbexample}{def:p18-is-dynamic-app-boundaries}{Two explicit applicability failures}
For the first distinction take \(r:\Rightarrow p\) and
\(s:\mathord{\sim}p\Rightarrow\neg\operatorname{app}(r)\).
Under priority-free direct exclusion, \([r]\) undercuts \([s]\),
blocking the reverse exclusion defeat; \(p\) is justified.
If \(r\) is instead replaced by
\(r':\mathord{\sim}\neg\operatorname{app}(r)\Rightarrow p\),
the two arguments mutually undercut and neither is justified.

For the second failure take
\[
 r:\Rightarrow p,\qquad
 s:\Rightarrow\neg\operatorname{app}(r),\qquad
 k:\Rightarrow s\triangleleft r.
\]
The support \(X=\{[s],[k]\}\) is conflict-free. Its priority suppresses
the sensitive exclusion of \([r]\), so both \([r]\) and \([s]\)
belong to \(G(X)\). But \([s]\) exclusion-attacks \([r]\).
If ranked applicability is desired within the proved core, a
separately specified positive-premise system may use
\[
 \rho_r:\Rightarrow\operatorname{app}(r),\quad
 r:\operatorname{app}(r)\Rightarrow p,\quad
 s:\Rightarrow\neg\operatorname{app}(r),\quad
 k:\Rightarrow s\triangleleft\rho_r.
\]
Here the priority supports the applicability default and defeats
the opposing inapplicability argument. This is a different theory,
not an equivalence claim for the failing direct variant.
\end{hbexample}

\section{Implementation and the next source obligations}

For a finite set of ground rules, enumerate legal argument contents,
compute their strict closures and finite frontier possibilities,
then compute \(G\) from the empty set. An increasing chain of subsets
of a universe of \(N\) arguments stabilizes after at most \(N\)
strict increases. The certificate proof constructs explanations
with decreasing finite stage ranks. A finite collection of schemata
with function symbols need not have a finite ground instantiation.
Finitariness likewise gives no effective way to verify that every
attacker has been found. The source reports an implementation based
on logic-programming techniques and discusses partial execution;
that report is not evidence of an implementation or performance
measurement for the present manuscript.

The dynamic operator and support-explicit procedure are
established above. Chapter~\ref{sec:p18-is-dynamic-completeness}
resolves the source single-response completeness claim by
counterexample and proves the corrected sequential dialogue.
Its attack-witness protocol has the exact dynamic overruling
correspondence of Theorem~\ref{thm:p18-is-licensing-reversed};
this is a specified protocol, not a bare swap of the old players.
A consistent replacement for priority-sensitive direct exclusion
remains a separate semantic obligation.
The earlier fixed-defeat LIS bridge applies when a fixed relation
has actually been supplied; it cannot freeze \(D_X\) during this
changing-support construction.
The source's unrestricted semi-decidability assertions also require
their own coding and scope proof before becoming algorithmic claims
of this reference.
